\documentclass[a4paper,12pt]{article}
\usepackage{graphicx, setspace, array, xcolor, mathtools, contour, tikz, amsthm, setspace, amsmath,amsfonts,amssymb, subcaption, fancyhdr, lipsum,multicol, geometry, gensymb}
\usepackage{colortbl}
\geometry{a4paper, margin=0.9in}
\contourlength{0.5pt}
\usetikzlibrary{patterns}
\usepackage[english]{babel}
\renewcommand{\qedsymbol}{$\blacksquare$}
\definecolor{darkgreen}{rgb}{0.0, 0.5, 0.0}
\pagestyle{fancy}
\fancyhf{}
\renewcommand{\headrulewidth}{0.4pt}
\fancyhead[L]{\rightmark}
\fancyhead[R]{\thepage}
\title{Modern Algebra}
\author{Assignment 4\\ \\ Yousef A. Abood\\ \\ ID: 900248250}
\date{October 2025}
\setlength{\parindent}{0pt}
\singlespacing
\parskip=1mm
\setcounter{section}{3}
\setcounter{subsection}{1}
\onehalfspacing
\begin{document}
\maketitle
\noindent\makebox[\linewidth]{\rule{15cm}{0.4pt}}
\section{Cyclic Groups}
\subsection*{Problem 1}
    By Corollary 4.2.8., we know that the generators of the group $(\mathbb{Z}_n,+)$ are the elements of the set of units $U_n.$ Thus,
    \begin{align*}
        &\text{the set of generators of } \mathbb{Z}_6 = U_6=\{1,5\} \\
        &\text{the set of generators of } \mathbb{Z}_8 = U_8=\{1,3,5,7\}\\
        &\text{the set of generators of } \mathbb{Z}_{20} = U_{20}=\{1,3,7,9,11,13,17,19\}.
    \end{align*}
\subsection*{Problem 5}
Since $3^3=7$, the order of $\langle 3\rangle=4$ and $\gcd(3,4)=1$, then $\langle 3\rangle=\langle 7\rangle$:
\begin{align*}
    &\langle 3\rangle=\{1,3,7,9\}\\
    &\langle 7\rangle=\{1,3,7,9\}.
\end{align*}
\subsection*{Problem 10}
Pick a generator element $g$ from the group $G=\mathbb{Z}_{24}$, the order of $G$ is $n=24$. Every generator of order $8$ must generate a cyclic subgroup of order $8$, by Corollary 4.2.3..
Choose any $k\in \mathbb{Z} \wedge 0 < k<n$, we know that any element in $G$ can be represented by $g^k$. We need $|g^k|=8$, by Theorem 4.2.5., we have:
\begin{align*}
    |g^k|= 8= \frac{n}{\gcd(n,k)}=\frac{24}{\gcd(24,k)} \implies \gcd(24,k)=3.
\end{align*}
Thus, all values of $k=21,15,9,3$, and the generators are the elements of the set $\{g^3,g^9,g^{15},g^{21}\}$, where $g$ is any generator of the group. We know that the generator of $(\mathbb{Z}_{24},+)=U_{24}=\{1,5,7,11,13,17,19,23\}.$ Therefore, $g \in U_{24}.$
\subsection*{Problem 11}
\begin{proof}
Let $G$ be a group and let $a\in G$. We want to prove that $\langle a \rangle =\langle a^{-1} \rangle.$ We need to show that $\langle a^{-1} \rangle \subseteq \langle a \rangle \wedge \langle a \rangle \subseteq \langle a^{-1} \rangle.$ First, we know that $\langle a \rangle$ is the set contains every $a^k$, where $k \in \mathbb{Z}.$ We observe that $(a^{-1})^k=(a^k)^{-1}=a^{-k},$ clearly belongs to $\langle a \rangle$ for any choice of $k.$ Hence, we showed that $\langle a^{-1} \rangle \subseteq \langle a \rangle.$ To prove that $\langle a \rangle \subseteq \langle a^{-1} \rangle$ we observe that for any $a^k \in \langle a \rangle$, we can write it as $(a^{-1})^{-k},$ and we know that $-k$ is an integer, so it is clear that $\langle a \rangle \subseteq \langle a^{-1} \rangle$ . Therefore, $\langle a^{-1} \rangle = \langle a \rangle$  
\end{proof}
\subsection*{Problem 27}
\begin{proof}
    We use the Euler formula to represent the complex numbers, that is, $z=e^{\frac{2\pi i}{n}}, z \in \mathbb{C}^*$. We need to show that each element has an order n. Observe that $(e^{\frac{2\pi i}{n}})^n=e^{2\pi i}=1.$ Next, pick an element $1\le k<n, k\in \mathbb{Z},$ and assume that $(z^k=e^{\frac{2\pi i}{n}})^k=e^{\frac{2k\pi i}{n}}=1.$ Hence, we need the power to be an integer multiple of $2 \pi$, so we need $\frac{k}{n}$ to be an integer. Thus, $k$ is a multiple of $n$, but that contradicts with our assumption that $k<n.$ Therefore, $n$ is the least integer such that $z^n=1$ and $n$ is the order of the cyclic group generated by any $z.$
\end{proof}
\subsection*{Problem 30}
\begin{proof}
    Suppose $G$ is a group with more than one element, so the group is not the trivial group. Since the group is closed under the operation, we choose element $g \in G \wedge g \neq e$ that generates $\langle g \rangle$. Since $\langle g \rangle$ is a subgroup and we know that the only available subgroups are the trivial one and the group itself. Thus, $\langle g \rangle=G$, so $G$ is cyclic. For the next part, by Corollary 4.2.6. and Corollary 4.2.3., we know that in finite groups, the order of any element in the group devides the order of the group and the order of any element is the order of the subgroups it generates. Hence, if the group is of a prime order $p$, then it does not have any subgroup of order that devides $p$ except $1,p$. We already know that the group has only two subgroups, which are the trivial group and the group itself, wich are of order $1,p$ respectively. So it is clear the group is of prime order.
\end{proof}
\subsection*{Problem 36}
$$
\begin{tabular}{|c|c|c|c|c|}
    \hline
    & $\mathbf{4}$ & $\mathbf{8}$ & $\mathbf{12}$ & $\mathbf{16}$ \\
    \hline
    $\mathbf{4}$ & 16 & 12 & 8 & 4 \\
    \hline
    $\mathbf{8}$ & 12 & 4 & 16 & 8 \\
    \hline
    $\mathbf{12}$ & 8 & 16 & 4 & 12 \\
    \hline
    $\mathbf{16}$ & 4 & 8 & 12 & 16\\
    \hline
\end{tabular}
$$
We know that the multiplication modulo 20 is associative. We observe the cayley table above and see that the group is closed under the operation as all the products belongs to the group. The identity element exists and it is $16.$ The inverses also exists, as for each number $a$ in the set there is a number $b$ in the set such that $ab=16.$ We observe that :$\{8^1, 8^2, 8^3, 8^4\}=\{8, 4,12,16 \}$ which is the whole group, so $8$ is a generator of the group and this group is cyclic. We also find that $\langle 12 \rangle =\{12,4,8,16\}.$ Thus, the generators of that group are $8,12.$
\subsection*{Problem 48}
If an element is divisible by $8$ and $10$, then it is divisible by its thier least common multiple, that is, $lcm (8,10)=40.$ We can now write $H=\{40k| k \in \mathbb{Z}\}.$ We show that $H$ is a subgroup of $\mathbb{Z}.$ It is clearly closed under addition, as $(40k+40n)$ is clearly divisible by $40$, where $n \in \mathbb{Z}$ . For the inverse, observe that $-(40k)$ is the inverse and it exitst in the set. If $k=0$, then $(40k)=(40\times 0)=0$ which is the identity element. If the property of the group changed to for any $n\in \mathbb{Z}$, $n$ is divisible by $8$ or $10$, then $(H, +)$ does not construct a subgroup. A counterexample is $8,10$ which are divisible by $8,10$ respectively, observe their additon: $8+10=18$ is not in the set as it does not satisfy the property. Thus, $H$ is not closed under addition.   
\subsection*{Problem 59}
\begin{proof}
    Suppose $G$ is a group. We want to proof that for any group there cannot be exactly two elements of order 2. For the sake of contradiction, we assume that there are exactly two elements with order 2. Choose $a,b \in G$ where $a,b$ are elements of order 2 and they are not the identity. If an element is of order 2, then it is the inverse of itself. By our choice, we know that $a \ne b \ne e.$ Observe the element $aba$, it has to be in the group. If $aba=a$, then $ab=e$ and $a=b^{-1}=b, b=a^{-1}=a$, which contradicts with our choice, so $aba \ne a$.  When $aba=e$, we see that $ba=a$ and $b=e$, which also contradicts with our choice. It cannot also be $b$ as
    \begin{align*}
        aba=b \implies
        ba=a^{-1}b=ab.
    \end{align*} 
We showed it implies commutativity if $aba=b.$ We take the square of $ab:$ 
\begin{align*}
    (ab)^2=a^2b^2=e\cdot e=e.
\end{align*}
Since $ab \ne b \ne a \ne e$, otherwirse $a=e, b=e, a=b$ which are contradiction with our choice. Thus, we showed that $ab$ is a distict element with order $2$ if $a,b$ commute, a contradiction with our assumption. If $a,b$ do not commute, then take the square of $aba$: 
\begin{align*}
    (aba)^2=abaaba=abeba=abba=aea=aa=e.
\end{align*} We also showed another distict integer that has order 2, a contradiction. Threfore, we proved that no group can have exactly two elements of order 2.
\end{proof}
\subsection*{Problem 61}
\begin{proof}
    Choose an arbitrary element $m\in H$, where $H$ is $\langle a \rangle \cap \langle b \rangle$. By Corollary 4.2.6., the order of element in a finite cyclic group must divide the order of the group, that is, the order of $m$ must divide $|\langle a \rangle|$ and $|\langle b \rangle|$. But we know that the only common divisor of $10$ and $21$ is $1$. The only element that has the order 1 is the identity element. Therefore, $\langle a \rangle \cap \langle b \rangle=\{e\}.$ 
\end{proof}
\end{document}